% ---------------------------------------------------------------------
%  Chapitre 5 — Le protocole, ses limites, ses suites
%  Rôle : livrer L'OUTIL (le produit du stage), assumer ses limites,
%  et tracer les étapes suivantes (dont les chantiers retirés du
%  périmètre : logistique sur grappes, signaux d'information, lead-lag).
% ---------------------------------------------------------------------
\chapter{Le protocole, ses limites, ses suites}
\label{chap:protocole}

\section{Le protocole récapitulé}
\label{sec:protocole}

% [À RÉDIGER — présenter l'outil comme le LIVRABLE du stage, 5-8 lignes]
% Idées :
%  - Ce que le stage livre n'est pas une stratégie gagnante mais une
%    procédure de jugement, réutilisable telle quelle sur tout signal
%    futur de l'agent (technique, informationnel ou relationnel).
%  - Chaque étage répare une condition d'application précise ; le
%    verdict n'est prononcé que si TOUS les étages concordent.
%  - Renvoyer au tableau récapitulatif.

\begin{table}[H]
\centering
\caption{Le protocole de validation en six étages~: chaque étage répare
une menace identifiée, et le verdict exige de tous les franchir.}
\label{tab:protocole}
\small
\begin{tabular}{@{}p{0.30\linewidth}p{0.34\linewidth}p{0.26\linewidth}@{}}
\toprule
Étage & Menace traitée & Outil \\
\midrule
1. \edgevar{} apparié au hasard & dérive du marché & Monte Carlo apparié \\
2. Student unilatéral + TCL & variabilité d'échantillonnage & test de référence \\
3. Seuil de Bonferroni & multiplicité (10 tests) & contrôle du risque global \\
4. DEFF + bootstrap par grappes & dépendance intra-titre & sondages, rééchantillonnage \\
5. Les deux portes (titre, mois) & pseudo-réplication & agrégation équipondérée \\
6. Série mensuelle AR($p$) & autocorrélation résiduelle & variance de long terme \\
\bottomrule
\end{tabular}
\end{table}

\section{Limites assumées}
\label{sec:limites}

% [À RÉDIGER — commentaire de l'encadré, 3-5 lignes] Idées :
%  - Aucune de ces limites n'invalide le verdict ; chacune trace en
%    revanche la frontière de ce que le protocole PEUT affirmer.
\begin{limites}
\begin{itemize}
  \item Résultats établis \emph{in-sample}~: le protocole juge des
        signaux passés~; une validation hors échantillon
        (\emph{walk-forward}) reste nécessaire avant tout usage
        opérationnel.
  \item Périmètre du jugement~: dix stratégies déclarées \emph{a
        priori}. À plus grande échelle (des milliers de variantes
        essayées), le contrôle du \emph{data snooping} exigerait des
        outils dédiés comme le \emph{Reality Check} de
        White~\cite{white2000}.
  \item Frais et conditions d'exécution modélisés de façon simplifiée.
  \item Puissance statistique élevée sur grands échantillons~:
        distinguer soigneusement significativité et taille d'effet.
\end{itemize}
\end{limites}

\section{Perspectives}
\label{sec:perspectives}

% [À RÉDIGER — 3 paragraphes courts, un par perspective] Idées :
%  1. VALIDATION TEMPORELLE : figer les seuils sur une période
%     d'apprentissage, juger sur la suivante (walk-forward /
%     TimeSeriesSplit) — la validation croisée classique est interdite
%     ici (fuite du futur).
%  2. L'EDGE CONDITIONNEL : modéliser P(win | contexte) par une
%     régression logistique. Écartée de ce mémoire PAR RIGUEUR : sur
%     38 210 trades dépendants, les p-values des coefficients seraient
%     trop optimistes (même piège qu'au chap. 3). Les outils adaptés
%     (modèles à effets mixtes, équations d'estimation généralisées)
%     dépassent le cadre du M1 -> prochaine étape naturelle.
%  3. LES AUTRES SIGNAUX DE L'AGENT : l'agent collecte aussi des
%     signaux d'information (contrats publics, transactions du
%     Congrès, réglementation) et des relations d'avance/retard entre
%     titres (causalité de Granger \cite{granger1969}). Des travaux
%     exploratoires ont été engagés sur ces deux fronts ; le protocole
%     du chap. 3 est précisément l'outil qui permettra de les juger —
%     hors du périmètre de ce mémoire.